% !TEX root = ../diary-book.tex

\section{令和7年11月1日 中間報告発表に向けて}

シミュレーションの結果から引き出せることを精一杯引き出す。
\begin{itemize}
  \item ミューオンのエネルギーごとの地中・海中への到達レート
  \item そのようなミューオンと一次粒子エネルギーの相関関係。もし 10 - 1000 TeV 領域の一次粒子、あるいはそれ以上の粒子に対する感度が高いことが分かれば、シャワー面のほうが優位になる可能性がたかい。地中にとどかないミューオンを活用出来る可能性が高まるため。ミューオンの半径は小さい？大きい？
  \item 結局角度依存性はまだシミュレーションできていないので、これについては早急にシミュレーションを回す必要があるだろう。
  \item どのくらいの領域までの一次宇宙線が影響するかについては半径でシミュレーションをしたとおりである。一次宇宙線のエネルギーが 100 GeV 以上では実はミューオンは電磁成分の半径の２倍くらい遠くまで届いているので、外側からの寄与という意味では４倍程度大きい事が想定される。ただしエネルギーは考えていない。電磁成分のほうが早く落ちるから半径が小さいのだとは思うけれど。
  \item 最大の繰り返しレートを決定することがシステムの全体像を規定する。極端な話地表がすべて検出器で埋められた際には、地中でのミューオンの測定レートが最大の繰り返しレートとなり、そこに検出効率がかかってくるという話になる。
\end{itemize}

\subsection{平面近似によるシャワー面解析}

明野のデータをつかってシャワー面の解析を行う。4つ以上のヒットを用いた慣性主軸の導出をおこなって、シャワー面の傾き（面の法線ベクトル）とシャワー面の位置を求める。


\subsection{平面近似}

\section{令和７年11月2日}

\subsection{ミューオンの深度}

ミューオンのエネルギー損失はほぼ 2 MeV/(g/cm$^{2}$) と思って良い。密度が $\rho$ g/cm$^{3}$ 程度の物体を通過すると考える。このとき $T$~MeV の運動エネルギーを持つミューオンが到達できる距離$L$は、ブラッグピーク付近のエネルギー損失を忘れると、$ T = 2 * L * \rho$ の関係があるため、$L = T / 2\rho$~cm と表される。水 500~m を通過できるミューオンのエネルギーは $T = 2 * 50000 * 1 = 100$~GeV 程度のものとなる。このくらいのエネルギーのものが飛来する頻度はどうなるか？

\subsection{ミューオンの到来頻度について}

ミューオンの到来頻度はシミュレーションを用いてある程度わかってきている。 100~TeV までの一次宇宙線を考慮すると 0.0014~個/sr/m$^{2}$/s が飛来することになる。ミューオンの分布の半径は 600 m 程度だとすると面積$1.2\times 10^{6}$~m$^{2$}の領域の一次宇宙線が関与できることとなる。また、立体角は45度くらいの天頂角までが重要であると仮定すると、\pi*(1 - $\sqrt{2}/2$) = 0.3 $\pi$ ~ 1.くらいの立体角になる。するとミューオンの到来数は 100~m$^{2}$の検出器にたいして 1600 個程度。これは 1~m の検出器にたいして 16 個程度であり、手のひらに１秒一個と比べると５倍ほど違うが、的外れというわけでもない。


\section{令和7年11月6日 シミュレーション}

過去の結果がどこにあるかを紐解いていたが、ata03 に存在することがわかった。データのカタログを作らないといけない。
kpro-main1 でやっているシミュレーションではシャワー面の同定効率が 100~TeV 領域でも100\%にはならず、それ以下のエネルギー領域でもミューオン単体の収量のほうが多いという結果が得られている。違いを分析する必要がある。

%論文を書くまではいかなくとも、どの程度の精度で測位が可能かを示すことができるとよい。
%また、応用先についての具体的な検討が必要であるので、それを含めて、どの程度までいけるかを整理する。
%その他の測位手法との差異についても整理を行う必要がある。
%
%\section{令和7年10月22日 ユースケースの一つとしての潜水艦の測位}
%
%これまでにいただいた情報を整理すると、潜水艦の測位のユースケースとしては以下のような条件が想定される。
%
%\begin{itemize}
%  \item 深さ：～500m
%  \item 速度：6km/h
%  \item 範囲：1km四方
%\end{itemize} 
%\subsection{深さとミューオンの到達可能性}
%深さ 500 m の場合、海水の密度は 1.03 g/cm$^3$ 程度であるので、500 m の水深では 515 m の水の厚みがあることになる。空気中の大気圧は 1 atm = 1013 hPa = 1013 mbar = 1013 g/cm$^2$ であるので、500 m の水深では 515 g/cm$^2$ の厚みがあることになる。ミューオンのエネルギー損失は 1 g/cm$^2$ あたり 2 MeV 程度であるので、500 m の水深では 1 GeV 程度のエネルギー損失があることになる。したがって、500 m の水深においても、1 GeV 以上のエネルギーを持つミューオンは到達することができる。ただし、1~GeV はギリギリのエネルギーであるので、2 GeV 以上のエネルギーを持つミューオンが到達することができると考えたほうがよい。また、測位には角度がついたミューオンが必要である。４５度程度までのミューオンを考えるとするとエネルギー損失は 1.4~GeV であり、2~GeV 以上のエネルギーを持つミューオンの到達が考えられる。
%
%\begin{codebox}{連絡先}
%  \begin{minted}{text}
%***************************************
%株式会社IHI
%航空・宇宙・防衛事業領域　防衛事業連携統括部
%高橋　慶貴
%TEL：03-6204-7120
%MAIL：takahashi0046@ihi-g.com
%***************************************
%\end{minted}
%\end{codebox}
%
%
%
